\chapter{INTRODUCTION}

Managing college events is typically a complex process involving multiple stakeholders, manual approvals, and dispersed communication channels. The College Event Management System is a comprehensive, full-stack web application designed to digitalize and streamline this entire workflow. Developed using the Laravel 11 framework and adhering to the Model-View-Controller (MVC) architectural pattern, the system provides a robust, centralized platform for students, faculty, and administrators to seamlessly coordinate campus activities.
\\\\
To ensure security and proper delegation of responsibilities, the application integrates a strict Role-Based Access Control (RBAC) mechanism. Users are categorized into Super Admin, Faculty/Admin, and Student roles, each provided with an interactive, role-specific dashboard. Administrators have the capability to create, categorize, and monitor events through analytics, as well as approve or reject proposals submitted by students. Meanwhile, students benefit from a personalized portal where they can discover upcoming activities, track their registrations, and propose new events.
\\\\
The frontend of the system is engineered using Tailwind CSS, delivering a professional, modern, and highly responsive user interface. Features such as participant capacity limits, automated registration workflows, and taxonomy-based event filtering are integrated to solve the logistical challenges of campus activities. Ultimately, the College Event Management System replaces traditional, paper-based processes with an efficient, scalable, and user-centric digital solution that fosters greater engagement within the academic community.
\\\\

